\chapter*{คำนำ}
\addcontentsline{toc}{chapter}{คำนำ}

เอกสารคำสอนและคู่มือปฏิบัติการเรื่อง \textbf{คู่มือพัฒนาโดรนเกษตรอัจฉริยะ (สถาปัตยกรรมแอปพลิเคชันและการวิเคราะห์ภาพถ่ายเชิงปริมาณ)} จัดทำขึ้นเพื่อใช้ประกอบการเรียนการสอนและการวิจัยในสาขาวิชาฟิสิกส์ประยุกต์ วิทยาการคอมพิวเตอร์ และเทคโนโลยีการเกษตรดิจิทัล มหาวิทยาลัยราชภัฏรำไพพรรณี โดยมีเป้าหมายเพื่อถ่ายทอดองค์ความรู้เชิงบูรณาการระหว่างฟิสิกส์ทัศนศาสตร์การสำรวจระยะไกล (Remote Sensing) ระบบสมองกลฝังตัวและการสื่อสารไร้สาย (Wireless IoT) และการพัฒนาแอปพลิเคชันบนสมาร์ทโฟนด้วยโครงร่างงาน Flutter

เนื้อหาภายในเล่มครอบคลุมองค์ความรู้ตั้งแต่ระดับพื้นฐานจนถึงการประยุกต์ใช้งานจริงในแปลงเกษตร เริ่มต้นจากสถาปัตยกรรมฮาร์ดแวร์ของโดรน โครงสร้างโปรโตคอลการสื่อสารข้อมูลแบบ UDP Datagram Socket ความถี่ 20 ครั้งต่อวินาที การแปลงสัญญาณคำสั่งการบินเป็นรหัสไบต์ฐานสิบหก การรับสตรีมภาพความเร็วสูงผ่านเครือข่ายไร้สาย และการออกแบบส่วนต่อประสานผู้ใช้สไตล์ Cockpit Heads-Up Display พร้อมการซ้อนภาพข้อมูลโทรมาตรเสมือนจริง (AR Telemetry) นอกจากนี้ยังนำเสนอระบบการบินอัตโนมัติตามพิกัดจุดบิน (Autonomous Waypoint Survey) และการคำนวณขอบเขตแปลงเกษตรกรรม

จุดเด่นสำคัญของคู่มือเล่มนี้คือการวิเคราะห์ภาพถ่ายเชิงปริมาณ โดยนำเสนออัลกอริทึมการคำนวณดัชนีพืชพรรณในช่วงคลื่นแสงที่ตามองเห็น ได้แก่ VARI, GLI, ExG และ Visible NDVI พร้อมการประมวลผลคู่ขนานแบบ Pure Dart Isolate เพื่อสร้างแผนที่ความร้อนแสดงระดับความสมบูรณ์ของคลอโรฟิลล์ การตัดแยกพื้นที่ทรงพุ่มพืช การประเมินขนาดพื้นที่จริงเป็นตารางเมตรและหน่วยไร่ไทย การตรวจนับจำนวนต้นพืชและระบุพิกัดกึ่งกลางทรงพุ่ม การตรวจจับจุดเว้าแหว่งเพื่อวางแผนปลูกซ่อม ตลอดจนการคำนวณอัตราการพ่นสารชีวภัณฑ์และประสิทธิภาพการใช้พลังงานของอากาศยานไร้คนขับ

คณะผู้ออกแบบและพัฒนาหวังเป็นอย่างยิ่งว่าเอกสารคำสอนและคู่มือปฏิบัติการเล่มนี้ จะเป็นประโยชน์แก่นักศึกษา คณาจารย์ นักวิจัย ตลอดจนเกษตรกรและผู้สนใจทั่วไป ในการนำเทคโนโลยีอากาศยานไร้คนขับและปัญญาประดิษฐ์ไปยกระดับการเกษตรแม่นยำของประเทศไทยให้ก้าวหน้าอย่างยั่งยืนสืบไป

\vspace{18mm}

\begin{flushright}
    \textbf{ผู้ช่วยศาสตราจารย์ ดร.ชีวะ ทัศนา}\par
    \textbf{ผู้ช่วยศาสตราจารย์ ดร.จิรภัทร จันทมาลี}\par
    \vspace{4pt}
    สาขาวิชาฟิสิกส์ คณะวิทยาศาสตร์และเทคโนโลยี\par
    มหาวิทยาลัยราชภัฏรำไพพรรณี\par
    ตุลาคม 2569
\end{flushright}

\clearpage
